\section{Setup}
\label{sec:protocol}

\paragraph{Benchmark and judged lists.}
We use \svo{}~\cite{hendricks2021probing} in two roles. As the ITM benchmark, it has 33,685 triplets of a caption, a positive image and a negative image, and a decision is scored on the two pairs of each triplet. As a retrieval collection, its 13,285 images are the candidates of a caption query (\ttoi{}) and its 10,978 captions those of an image query (\itot{}); the \emph{strict} Success@1 counts only a benchmark positive of the query as correct. The judged lists cover 100 random caption queries and 100 random image queries, with the top 10 candidates of five embedding models for each: CLIP ViT-L/14~\cite{radford2021learning}, the BLIP-2 Q-Former~\cite{li2023blip2}, FLAVA~\cite{singh2022flava}, SigLIP2~\cite{tschannen2025siglip2} and \qemb{}~\cite{qwen3vlembedding2026}. The order in which a model returned its ten candidates, by their similarity with the query, is the \emph{original order} of the list. Every list was labelled blind to its model, the image-query lists of the first four models by three annotators with a majority vote (Fleiss' $\kappa$ 0.64) and the other lists by one expert; a pair is correct when the caption describes the image fully, and an incorrect pair carries the wrong role (subject, verb or object).

\paragraph{Matching decisions.}
The \emph{BLIP head}~\cite{li2022blip} (ViT-B, the COCO retrieval checkpoint), the \emph{BLIP-2 head} of the same BLIP-2 model and the \emph{FLAVA head} are the image--text matching heads of these models; their match probability is the score and 0.5 the threshold. The \emph{Qwen2.5-VL prompt} asks the zero-shot Qwen2.5-VL-7B-Instruct model whether the image matches the caption in subject, verb and object; the score is the probability of ``yes'' against ``no'' at the first generated token, the threshold 0.5. The \emph{\qemb{} yes/no} decision embeds the pair under a classification instruction and scores it by the difference of its similarities with the embeddings of ``Yes'' and ``No'', the threshold 0. \emph{Qwen3-VL-Reranker}~\cite{qwen3vlembedding2026} is a released multimodal re-ranker (2B parameters) with its default instruction; its score is the same probability of ``yes'' against ``no'', the threshold 0.5. \emph{Grounding CLIP}~\cite{vongala2025compositional} is a training-free compositional score: it adds the CLIP ViT-L/14 embeddings of the subject and object regions found by Grounding DINO, weighted by their similarity with their phrases, to the whole-image embedding before the cosine with the caption; it has no threshold, so its benchmark number is the pairwise accuracy. Per pair, the heads cost a 0.2--1.1B-parameter forward pass, the two Qwen3 decisions a 2B one and the prompt a 7B one.

\paragraph{Measures.}
We report the AUC of each matching decision's score over all judged pairs, and the share of judged-correct pairs its default threshold accepts (TPR) and of judged-incorrect pairs it rejects (TNR). Inside a top-10 list we compare the AUC of its score with that of the original order. We then use each decision on the lists in two ways. As a re-ranker, it re-orders the ten candidates by its score, with ties kept in the original order. As a filter, it moves the candidates it accepts in front of those it rejects. We measure the judged Success@1 before and after; the first candidate counts when it is judged correct, whether it is the benchmark positive or not. Two further variants are tuned on the judged lists of the other queries in ten random folds: a filter whose threshold maximises the judged Success@1 of the training lists, and a re-ordering by the score interpolated with the embedding model's similarity (both standardised inside the list), with the weight tuned likewise. Every change is tested with a paired randomisation test on the per-query differences~\cite{smucker2007comparison}, with Holm's correction~\cite{holm1979simple} over the ten list-and-direction cells of a decision for the per-list changes and over the fourteen decision-and-direction cells for the mean changes.
